% report/main.tex
\documentclass[10pt,a4paper]{article}

% --- Language / fonts (pdfLaTeX-friendly) ---
\usepackage[T2A]{fontenc}
\usepackage[utf8]{inputenc}
\usepackage[russian,english]{babel}

% --- Layout ---
\usepackage[a4paper,margin=1.7cm]{geometry}
\usepackage{microtype}
\setlength{\parindent}{0pt}
\setlength{\parskip}{0.35em}
\usepackage{enumitem}
\setlist[itemize]{noitemsep,topsep=0.2em,leftmargin=1.25em}
\setlist[enumerate]{noitemsep,topsep=0.2em,leftmargin=1.25em}

% --- Figures / tables ---
\usepackage{graphicx}
\usepackage{caption}
\usepackage{subcaption}
\usepackage{booktabs}
\usepackage{float}
\usepackage{multirow}
\usepackage{xcolor}

% --- Links / refs ---
\usepackage[hidelinks]{hyperref}

% --- Code / verbatim inputs ---
\usepackage{fancyvrb}

% --- Bibliography compact ---
\usepackage[numbers]{natbib}
\renewcommand{\bibsection}{\section*{References}}
\setlength{\bibsep}{0pt}

% Paths
\graphicspath{{figures/}}

% ---------- Macros ----------
\newcommand{\method}[1]{\textbf{#1}}
\newcommand{\tok}[1]{\texttt{#1}}
\newcommand{\placeholder}[1]{\textcolor{gray}{\textit{#1}}}

\title{\vspace{-0.8em}\textbf{Fine-tuning NanoVLM для управления агентом в MiniGrid EmptyEnv: SFT vs GRPO-action vs GRPO-text+action}\vspace{-0.4em}}
\author{\placeholder{Автор: ФИО / группа}}
\date{\vspace{-0.8em}\placeholder{Дата: YYYY-MM-DD}}

\begin{document}
\maketitle
\vspace{-0.8em}

\section{Постановка}
Нужно адаптировать vision-and-language модель \method{NanoVLM} для управления агентом в среде \method{MiniGrid EmptyEnv} \citep{minigrid-emptyenv, nanovlm-repo}.
Сравниваются три режима дообучения:
\begin{itemize}
  \item \method{SFT}: supervised fine-tuning на экспертных траекториях.
  \item \method{GRPO-action}: RL-дообучение через Guided Regularized Policy Optimization, где модель \emph{напрямую} выдаёт действие \citep{grpo}.
  \item \method{GRPO-text+action}: GRPO, где модель генерирует короткий текст (2--3 предложения) + действие.
\end{itemize}

Метрики: \textbf{success rate} (доля эпизодов, где агент достиг цели в пределах horizon) и \textbf{return} (суммарная награда MiniGrid; детали см. документацию \citep{minigrid-emptyenv}). Цель сравнения: \textbf{(i)} итоговое качество, \textbf{(ii)} sample efficiency (env steps / episodes до достижения success=0.8 и 0.9), \textbf{(iii)} влияние формата вывода.

\section{Данные и эксперт}
\textbf{Среда.} EmptyEnv представляет собой пустую комнату с целью (goal); агент наблюдает RGB-изображение/визуальные токены и должен добраться до goal \citep{minigrid-emptyenv}.

\textbf{Эксперт.} Для генерации датасета используется детерминированный эксперт-планировщик на полном состоянии среды (например, shortest-path planning). В EmptyEnv отсутствие препятствий позволяет строить план как последовательность поворотов и шагов вперёд до достижения цели.
\emph{Почему так:} эксперт даёт близко-оптимальные траектории, прост в реализации и воспроизводим, а SFT становится чистым измерением «умеет ли модель имитировать поведение по визуальному наблюдению».

\textbf{Формирование датасета.} Записываются траектории $(o_t, a_t)$ (и опционально $r_t, done_t$) из $N$ эпизодов, где $o_t$ — наблюдение (изображение), $a_t$ — следующее действие эксперта.

\smallskip
\textbf{Примечание (логи/формат).} Скрипт генерации таблицы/графиков обычно ожидает логи в формате CSV/JSONL со столбцами (или ключами) минимум:
\tok{env_steps}, \tok{episode}, \tok{success} (0/1), \tok{return}, \tok{success_rate_rolling}, \tok{return_mean_rolling}.
Фигуры в \tok{report/figures/*.png} считаются построенными из этих логов.

\section{Методы}

\subsection{Форматы входа/выхода}
\textbf{Вход.} На шаге $t$ модель получает визуальное наблюдение $o_t$ и краткий текстовый промпт (например, «Выбери следующее действие»).

\textbf{Выход A: direct action token.}
Модель генерирует один токен действия из фиксированного словаря MiniGrid, например:
\[
a_t \in \{\tok{left},\tok{right},\tok{forward},\tok{pickup},\tok{drop},\tok{toggle},\tok{done}\}.
\]
(Точный словарь/маппинг на action-id соответствует реализации в коде.)

\textbf{Выход B: text (2--3 sentences) + action token.}
Модель генерирует короткое описание/план и затем действие в строгом формате:
\begin{quote}\small
\tok{PLAN: <sentence1> <sentence2> [<sentence3>]}\\
\tok{ACTION: <action\_token>}
\end{quote}
\emph{Обоснование:} короткий план может стабилизировать обучение в RL за счёт более структурированной внутренней репрезентации и уменьшения «случайных» действий; формат фиксирован (легче парсить и штрафовать отклонения).

\subsection{SFT (baseline)}
SFT минимизирует кросс-энтропию по следующему действию эксперта:
\[
\mathcal{L}_{\text{SFT}} = -\mathbb{E}_{(o_t,a_t)\sim \mathcal{D}} \log \pi_\theta(a_t \mid o_t, \text{prompt}).
\]
Для режима text+action в SFT (если он использовался как претрейн для GRPO-text+action) таргет может включать и токены плана, и \tok{ACTION}.

\subsection{GRPO-action}
Политика инициализируется весами SFT.
GRPO обновляет $\pi_\theta$ по группам сэмплов: для каждого состояния собирается группа из $G$ ответов (траекторий/действий), рассчитывается групповая нормализация преимущества и добавляется регуляризация к референс-политике $\pi_{\text{ref}}$ (часто это замороженная копия текущей/стартовой политики) \citep{grpo}.
Схематично оптимизируется:
\[
\max_\theta \ \mathbb{E}\left[\hat{A}\cdot \log \pi_\theta(a_t\mid \cdot)\right] - \beta \, \mathrm{KL}\!\left(\pi_\theta \,\|\, \pi_{\text{ref}}\right),
\]
где $\hat{A}$ — advantage на основе сравнения ответов внутри группы.

\subsection{GRPO-text+action}
Аналогично GRPO-action, но действие извлекается из строки \tok{ACTION: ...}. Reward считается по фактическому действию в среде, плюс возможны \textbf{форматные штрафы} (опционально): если нарушен шаблон, если текст слишком длинный, если отсутствует \tok{ACTION}, и т.п.
Это делает канал «планирования» контролируемым, но не даёт ему «захватить» оптимизацию.

\section{Результаты}
\subsection{Кривые обучения}
Рис.~\ref{fig:success} показывает success rate для трёх режимов. Возвраты (return) показаны на Рис.~\ref{fig:return}.
\begin{figure}[H]
  \centering
  \begin{subfigure}[t]{0.32\linewidth}
    \centering
    \includegraphics[width=\linewidth]{sft_success.png}
    \caption{SFT}
  \end{subfigure}
  \begin{subfigure}[t]{0.32\linewidth}
    \centering
    \includegraphics[width=\linewidth]{grpo_action_success.png}
    \caption{GRPO-action}
  \end{subfigure}
  \begin{subfigure}[t]{0.32\linewidth}
    \centering
    \includegraphics[width=\linewidth]{grpo_text_action_success.png}
    \caption{GRPO-text+action}
  \end{subfigure}
  \caption{Success rate по мере обучения (ось X: env steps или epochs — как в ваших логах/графиках).}
  \label{fig:success}
\end{figure}

\begin{figure}[H]
  \centering
  \IfFileExists{return_curves.png}{
    \includegraphics[width=0.98\linewidth]{return_curves.png}
  }{
    \placeholder{(return\_curves.png не найден: положите сюда общий график или замените на три отдельных)}
  }
  \caption{Return curves для трёх режимов (единый график или объединённые кривые).}
  \label{fig:return}
\end{figure}

\subsection{Сводная таблица и sample efficiency}
Таблица~\ref{tab:results} автоматически подставляется из \tok{report/tables/results_table.tex}.
\begin{table}[H]
  \centering
  \small
  \input{tables/results_table.tex}
  \caption{Сводные результаты: финальные метрики и sample efficiency. \placeholder{Числа подставляются из results\_table.tex.}}
  \label{tab:results}
\end{table}

\textbf{Sample efficiency сравнение.}
До порогов success=0.8 и success=0.9 сравниваются:
\begin{itemize}
  \item \textbf{env steps}: сколько шагов среды требуется до достижения порога (по сглаженному success rate).
  \item \textbf{episodes}: сколько эпизодов требуется до достижения порога.
\end{itemize}
\placeholder{В тексте отчёта не придумываем значения: см. Table~\ref{tab:results} (столбцы steps\_to\_0.8 / steps\_to\_0.9 и episodes\_to\_0.8 / episodes\_to\_0.9, если они присутствуют).}

\smallskip
\textbf{Примечание (ожидаемые поля таблицы).}
В \tok{results\_table.tex} обычно ожидаются строки для трёх методов и колонки вида:
\[
\text{final\_success}, \ \text{final\_return}, \
\text{steps\_to\_0.8}, \ \text{steps\_to\_0.9}, \
\text{episodes\_to\_0.8}, \ \text{episodes\_to\_0.9}
\]
(названия могут отличаться; таблица считается «истиной», а текст ссылается на неё).

\subsection{Гиперпараметры (из config.yaml)}
Ниже — конфиги запусков, подгружаемые из \tok{results/*/config.yaml}. Это нужно, чтобы отчёт был воспроизводимым (без копипасты в текст).
\begin{figure}[H]
\centering
\small
\begin{minipage}{0.32\linewidth}
\textbf{SFT config}
\IfFileExists{../results/sft/config.yaml}{
\VerbatimInput[fontsize=\scriptsize,frame=single]{../results/sft/config.yaml}
}{
\placeholder{(нет ../results/sft/config.yaml)}
}
\end{minipage}\hfill
\begin{minipage}{0.32\linewidth}
\textbf{GRPO-action config}
\IfFileExists{../results/grpo_action/config.yaml}{
\VerbatimInput[fontsize=\scriptsize,frame=single]{../results/grpo_action/config.yaml}
}{
\placeholder{(нет ../results/grpo\_action/config.yaml)}
}
\end{minipage}\hfill
\begin{minipage}{0.32\linewidth}
\textbf{GRPO-text+action config}
\IfFileExists{../results/grpo_text_action/config.yaml}{
\VerbatimInput[fontsize=\scriptsize,frame=single]{../results/grpo_text_action/config.yaml}
}{
\placeholder{(нет ../results/grpo\_text\_action/config.yaml)}
}
\end{minipage}
\caption{Гиперпараметры запусков (подставляются из YAML-конфигов).}
\end{figure}

\section{Короткое обсуждение и future work}
\textbf{Что смотреть в failure modes:}
\begin{itemize}
  \item \emph{Ошибки ориентации:} агент «крутится» из-за неустойчивого определения направления по визуальному наблюдению.
  \item \emph{Длинные эпизоды без прогресса:} разреженная награда и слабая exploration в RL.
  \item \emph{Форматные срывы в text+action:} модель забывает \tok{ACTION:} или начинает генерить слишком много текста (лечится ограничением длины/штрафами).
\end{itemize}

\textbf{Короткий итог.} \placeholder{Вывод формулируется по Table~\ref{tab:results}: кто быстрее достигает 0.8/0.9 success и у кого выше final success/return.}

\textbf{Future work (5--10 пунктов):}
\begin{itemize}
  \item Curriculum по размеру карты и max\_steps (от маленьких к большим).
  \item Балансировка датасета эксперта по направлениям/позициям цели (чтобы не доминировали «простые» случаи).
  \item Data augmentation для визуального входа (простые трансформации) + проверка устойчивости.
  \item PEFT/LoRA вместо полного fine-tune (дешевле и стабильнее), отдельно для vision и language частей.
  \item Явные форматные штрафы/награды для text+action (валидный шаблон, лимит токенов плана).
  \item Reward shaping: небольшая плотная компонента по приближению к цели (если допускается постановкой).
  \item Увеличение group size $G$ в GRPO и анализ trade-off compute vs sample efficiency.
  \item Анализ обобщения на другие EmptyEnv размеры/seed’ы (OOD по старту/goal).
\end{itemize}

\vspace{-0.2em}
\bibliographystyle{unsrtnat}
\bibliography{refs}

\end{document}